% =====================================================================
% Table --- the elasticities of the comparison (Sec. 4.2)
% ---------------------------------------------------------------------
% Requested by reviewer 1 (R1.3: "including elasticity values in a table").
% Values are read from the executed design experiments/designs/
% matrix_5x3_v11.toml and from the manifests of the reported cells
% (runs/matrix_5x3-v11-*): sigma, theta, epsilon, eta are the design values;
% eta_s is the bracket printed in Table 4; the two extension parameters are
% the ones the reported tables of Section 6 use. None of these is estimated
% from the input--output table.
% =====================================================================
\begin{table}[t]
\centering
\small
\begin{tabular}{@{}l p{5.9cm} l l@{}}
\toprule
Parameter & Meaning in the kernel & Value & Set in \\
\midrule
$\sigma$ & substitution across consumption goods & $0.9$ & all cells \\
$\theta$ & substitution across intermediates & $0.5$ & all cells \\
$\varepsilon$ & labor--intermediate nesting & $0.5$ & all cells \\
$\eta$ & intersectoral labor-allocation selector & $\{0,1\}$ & labor closure \\
$\eta_s$ & sectoral labor-supply elasticity with respect to the real wage & $\{0.5,1,2\}$ & elastic-supply column \\
\addlinespace

$\delta$ & capacity door: utilization elasticity of productivity & $\{0,0.05\}$ & Section~\ref{subsec:pricing-capacity} \\
$\kappa$ & markup door: margin--utilization slope & $\{0.25,0.5\}$ & Section~\ref{subsec:pricing-markup} \\

\bottomrule
\end{tabular}
\par\smallskip
\parbox{\textwidth}{\footnotesize
\textit{Note:} $\eta=0$ fixes employment in each sector at its baseline level,
with sector-specific wages; $\eta=1$ permits cost-minimizing labor
reallocation at a common wage.}
\caption{Parameter values for the baseline comparison and pricing extensions.}
\label{tab:parameters}
\end{table}
